% =======================================
% RÉSUMÉ (Français) + ABSTRACT (Anglais) 
% =======================================


% ----- Résumé français -----
\chapter*{Résumé}
\addcontentsline{toc}{chapter}{Résumé}
\markboth{Résumé}{Résumé}

Les attaques par déni de service distribué (\gls{ddos}) constituent l'une des menaces les plus critiques pour la disponibilité des services numériques modernes. Face à l'augmentation exponentielle de leur volume et de leur sophistication, les approches de mitigation centralisées montrent leurs limites~: points de défaillance uniques, coûts élevés et absence de coordination entre acteurs indépendants.

Ce mémoire a pour objectif de concevoir et d'implémenter un système distribué de mitigation collaborative des attaques \gls{ddos}, reposant sur une coalition pair-à-pair de n\oe{}uds de protection hétérogènes. Pour ce faire, le système intègre un mécanisme de confiance fondé sur le modèle PeerTrust avec crédibilité historique par session, un algorithme de sélection multicritère des pairs basé sur le Weighted Sum Model avec pondération \gls{ahp}, ainsi qu'un protocole de coordination P2P structuré en quatre phases permettant à un n\oe{}ud victime de solliciter, négocier et clore une session d'aide sans coordinateur central.

Le prototype a été développé en Node.js avec PostgreSQL et évalué sur une coalition simulée à différentes échelles. Les tests confirment la tenue en charge du système ainsi que la gestion correcte des cycles d'aide de bout en bout, depuis la détection de l'attaque jusqu'au recalcul des scores de confiance.

\medskip
\noindent\textbf{Mots-clés~:} DDoS, mitigation collaborative, pair-à-pair, confiance, PeerTrust, Weighted Sum Model, mTLS, JWT, sécurité réseau.

\newpage

% ----- Abstract anglais -----
\chapter*{Abstract}
\addcontentsline{toc}{chapter}{Abstract}
\markboth{Abstract}{Abstract}

Distributed Denial of Service (DDoS) attacks represent one of the most critical threats to the availability of modern digital services. As their volume and sophistication grow exponentially, centralized mitigation approaches show their limitations~: single points of failure, high costs, and lack of coordination between independent actors.

This thesis aims to design and implement a distributed collaborative DDoS mitigation system based on a peer-to-peer coalition of heterogeneous protection nodes. The system incorporates a trust mechanism based on the PeerTrust model with per-session historical credibility, a multi-criteria peer selection algorithm using the Weighted Sum Model with AHP weighting, and a P2P coordination protocol structured in four phases enabling a victim node to request, negotiate and close a help session without any central coordinator.

The prototype was developed in Node.js with PostgreSQL and evaluated on a simulated coalition at various scales. Tests confirm the system's ability to handle load and correctly manage end-to-end help cycles, from attack detection through trust score recalculation.

\medskip
\noindent\textbf{Keywords~:} DDoS, collaborative mitigation, peer-to-peer, trust, PeerTrust, Weighted Sum Model, mTLS, JWT, network security.

\newpage
